% ==========================================================================
%  ShinySwarm — REST API Architecture Benchmark Report
%  Generated from k6 results: 5 runs (run-1 .. run-5), Jun 2026
%  Cluster: Hetzner CX32 (8 vCPU / 16 GB RAM), K3s single-node
%  Stack: Plumber (R) + Redis + Spring Boot relay + PostgreSQL + Angular/nginx
%
%  Scope notes:
%   - Throughput (test 04) is EXCLUDED: runs are not independent (server not
%     reset between runs), so its numbers are invalid. See validity notes.
%   - Latency tests (01,02,03,08) ramp VUs but DO NOT tag vu_bucket, so their
%     avg/median/p95 mix all load levels. Reported with that caveat.
%   - All per-app numbers are the mean of the 5 runs (run-03 trial excluded).
% ==========================================================================
\documentclass[11pt,a4paper]{article}
\usepackage[margin=2.5cm]{geometry}
\usepackage{booktabs,tabularx,longtable}
\usepackage{xcolor,colortbl}
\usepackage{hyperref}
\usepackage{amsmath}
\usepackage[T1]{fontenc}
\usepackage{lmodern}

\definecolor{pass}{HTML}{2E7D32}
\definecolor{fail}{HTML}{C62828}
\definecolor{warn}{HTML}{F57F17}
\newcommand{\pass}{\textcolor{pass}{\textbf{PASS}}}
\newcommand{\fail}{\textcolor{fail}{\textbf{FAIL}}}
\newcommand{\warn}{\textcolor{warn}{\textbf{WARN}}}

\title{ShinySwarm Benchmark Report\\[4pt]\large REST API Architecture (5-run aggregate)}
\author{Iva Gunjaca — University of Amsterdam}
\date{June 2026}

\begin{document}
\maketitle

% ------------------------------------------------------------------
\section{Test Environment}
% ------------------------------------------------------------------

\begin{tabular}{@{}ll@{}}
\toprule
\textbf{Parameter}       & \textbf{Value} \\
\midrule
Server                    & Hetzner CX32 — 8 vCPU, 16\,GB RAM \\
Kubernetes                & K3s single-node \\
R backends                & Plumber APIs (HTTP) \\
State store               & Redis (live) + PostgreSQL (checkpoints) \\
Relay                     & Spring Boot API gateway \\
Frontend                  & Angular + nginx reverse proxy \\
Load generator            & k6 (remote) \\
Apps under test           & 8 collaborative R Shiny applications \\
Runs aggregated           & 5 (run-1 \ldots run-5; run-03 trial excluded) \\
\bottomrule
\end{tabular}

% ------------------------------------------------------------------
\section*{Validity Notes (read first)}
% ------------------------------------------------------------------
\addcontentsline{toc}{section}{Validity Notes}

\begin{itemize}
  \item \textbf{Throughput (test 04) is excluded} as invalid: the 5 runs were
        not independent (server not reset between runs), so its numbers are not
        reported here.
  \item \textbf{Latency tests lack a per-VU breakdown.} Tests 01, 02, 03 and
        08 ramp virtual users but tag metrics only with \texttt{app\_name},
        not \texttt{vu\_bucket}. Their avg/median/p95 therefore aggregate all
        load levels: \emph{the median reflects low-VU samples while the p95
        reflects the high-VU tail}. Read median as the typical case, p95 as the
        loaded tail. This matters more for REST than for Kafka because the REST
        tail is dominated by the synchronous relay under load (see
        \S\ref{sec:postpoll}).
  \item \textbf{Per-app thresholds are baseline-calibrated} (10$\times$ each
        app's single-user baseline). A PASS means the app stayed within its own
        envelope under load, not that it was fast in absolute terms.
\end{itemize}

% ------------------------------------------------------------------
\section{Summary}
% ------------------------------------------------------------------

\begin{tabularx}{\textwidth}{@{}lccX@{}}
\toprule
\textbf{\#} & \textbf{Benchmark} & \textbf{Result} & \textbf{Key Finding (5-run mean)} \\
\midrule
01 & State Relay Latency       & \warn & medians 21--90\,ms, but heavy p95 tails (2--8\,s); Advanced fails SLO; MLTrainer median 31\,s (intrinsic) \\
02 & Collaboration Latency     & \warn & medians 60--110\,ms; p95 2.6--7.6\,s; Advanced fails SLO \\
03 & Save \& Restore Latency   & \warn & median 10\,ms; p95 up to 1.1\,s; Geospatial, MonteCarlo, Advanced fail SLO \\
04 & Throughput                 & \multicolumn{2}{l}{\textbf{EXCLUDED — invalid (see validity notes)}} \\
05 & Data Loss                 & \pass & 0\% loss, all runs, all apps \\
06 & Cross-Contamination       & \pass & 0 contamination events, all runs \\
07 & Session Lifecycle          & \warn & 0 errors; op median 1.0\,s, p95 9.9\,s; Advanced fails SLO \\
08 & Multi-User Collaboration   & \pass & all apps 20--57\,ms; all within threshold \\
09 & WebSocket Presence         & \pass & connect $\approx$\,4\,ms median, stable across VUs \\
\bottomrule
\end{tabularx}

\medskip
\noindent\textbf{Headline.} At low load REST is fast for every app (relay
medians of tens of milliseconds). Under concurrency the p95 tails grow large
because the synchronous Spring Boot relay serialises requests: the cost shows
up in the POST leg, not the poll (see \S\ref{sec:postpoll}). MLTrainer is a
separate case — its $\approx$\,31\,s median is intrinsic R compute, not relay
overhead.

% ------------------------------------------------------------------
\section{Detailed Results}
% ------------------------------------------------------------------

% ---- 01 ----
\subsection{01 — State Relay Latency}

Round-trip: POST state $\to$ Spring Boot $\to$ Plumber computes $\to$ GET from
Redis. Per-app \texttt{relay\_total\_ms}, mean of 5 runs.

\begin{tabular}{@{}lrrrc@{}}
\toprule
\textbf{App} & \textbf{avg} & \textbf{median} & \textbf{p95} & \textbf{threshold} \\
\midrule
Calculator      & 1.2\,s  & 41\,ms  & 7.7\,s  & \pass \\
Geospatial      & 807\,ms & 21\,ms  & 4.9\,s  & \pass \\
DataExchange    & 491\,ms & 21\,ms  & 2.0\,s  & \pass \\
MonteCarlo      & 503\,ms & 88\,ms  & 2.2\,s  & \pass \\
Advanced        & 508\,ms & 30\,ms  & 2.0\,s  & \fail \\
Analytics       & 1.4\,s  & 46\,ms  & 8.0\,s  & \pass \\
ClimateAnomaly  & 1.4\,s  & 44\,ms  & 8.8\,s  & \pass \\
MLTrainer       & 38.7\,s & 30.8\,s & 78.0\,s & \pass\textsuperscript{*} \\
\bottomrule
\end{tabular}

\medskip
\noindent\footnotesize\textsuperscript{*}MLTrainer's threshold (143\,s) is
calibrated to its $\approx$\,14\,s single-user compute, so it passes despite a
31\,s median. Advanced fails its 730\,ms SLO due to its loaded p95.
\normalsize

\medskip
\noindent\textbf{Note the median/p95 gap.} Most apps have a median of tens of
ms but a p95 of several seconds. That gap is the VU ramp: low-VU samples are
fast, high-VU samples are slow, and they are pooled (no \texttt{vu\_bucket}
tag). MLTrainer is the exception — its median itself is large because the app's
computation is large.

% ---- post/poll split ----
\subsection{Where the time goes: POST vs.\ POLL}\label{sec:postpoll}

REST decomposes the relay into the POST (synchronous call through Spring Boot
to Plumber) and the POLL (read result from Redis). Global values (run-2,
representative):

\begin{tabular}{@{}lrr@{}}
\toprule
\textbf{Stage} & \textbf{median} & \textbf{p95} \\
\midrule
\texttt{relay\_post\_ms} (POST via Spring Boot) & 29\,ms & 30.0\,s \\
\texttt{relay\_poll\_ms} (read from Redis)      & 5\,ms  & 1.5\,s  \\
\bottomrule
\end{tabular}

\medskip
\noindent Unlike Kafka (where the wait-for-result poll dominates), in REST the
\textbf{POST leg} carries the tail: it is synchronous and runs through the
central Spring Boot relay, which serialises under load. The Redis poll stays
cheap. This is the architectural signature of the REST design — a single relay
point that becomes the bottleneck as concurrency rises.

% ---- 02 ----
\subsection{02 — Collaboration Latency}

Host posts state; guest polls until the update appears.
Per-app \texttt{collab\_propagation\_ms}, mean of 5 runs.

\begin{tabular}{@{}lrrrc@{}}
\toprule
\textbf{App} & \textbf{avg} & \textbf{median} & \textbf{p95} & \textbf{threshold} \\
\midrule
Calculator      & 640\,ms & 71\,ms  & 2.6\,s & \pass \\
Geospatial      & 915\,ms & 65\,ms  & 4.5\,s & \pass \\
DataExchange    & 757\,ms & 71\,ms  & 2.8\,s & \pass \\
MonteCarlo      & 818\,ms & 85\,ms  & 4.7\,s & \pass \\
Advanced        & 980\,ms & 107\,ms & 4.3\,s & \fail \\
Analytics       & 827\,ms & 97\,ms  & 5.0\,s & \pass \\
ClimateAnomaly  & 617\,ms & 60\,ms  & 3.1\,s & \pass \\
MLTrainer       & 1.6\,s  & 343\,ms & 7.6\,s & \pass\textsuperscript{*} \\
\bottomrule
\end{tabular}

\medskip
\noindent Median propagation is well under 350\,ms for every app; the
multi-second p95 is again the loaded tail.

% ---- 03 ----
\subsection{03 — Save \& Restore Latency}

PostgreSQL snapshot cycle. Per-app \texttt{save\_to\_db\_ms}, mean of 5 runs.

\begin{tabular}{@{}lrrrc@{}}
\toprule
\textbf{App} & \textbf{avg} & \textbf{median} & \textbf{p95} & \textbf{threshold} \\
\midrule
Calculator      & 251\,ms & 10\,ms & 590\,ms & \pass \\
Geospatial      & 365\,ms & 10\,ms & 717\,ms & \fail \\
DataExchange    & 251\,ms & 10\,ms & 1.1\,s  & \pass \\
MonteCarlo      & 304\,ms & 10\,ms & 844\,ms & \fail \\
Advanced        & 296\,ms & 10\,ms & 269\,ms & \fail \\
Analytics       & 271\,ms & 10\,ms & 407\,ms & \pass \\
ClimateAnomaly  & 132\,ms & 10\,ms & 47\,ms  & \pass \\
MLTrainer       & 265\,ms & 10\,ms & 201\,ms & \pass \\
\bottomrule
\end{tabular}

\medskip
\noindent Save is 10\,ms at the median for every app, but three apps cross
their (tight) save SLOs on the loaded tail. The PostgreSQL connection pool is
the shared contention point under concurrency.

% ---- 05 ----
\subsection{05 — Data Loss}

\begin{tabular}{@{}lr@{}}
\toprule
\textbf{Metric} & \textbf{Value (all 5 runs)} \\
\midrule
Loss rate     & 0\% \\
Checks passed & 105 / 105 (100\%) per run \\
\bottomrule
\end{tabular}

\medskip\noindent Threshold \texttt{dataloss\_pct $<$ 5\%}: \pass.

% ---- 06 ----
\subsection{06 — Cross-Contamination}

\begin{tabular}{@{}lr@{}}
\toprule
\textbf{Metric} & \textbf{Value (all 5 runs)} \\
\midrule
Contamination events & 0 \\
\bottomrule
\end{tabular}

\medskip\noindent Threshold \texttt{contamination\_events == 0}: \pass.

% ---- 07 ----
\subsection{07 — Session Lifecycle}

\begin{tabular}{@{}lr@{}}
\toprule
\textbf{Metric} & \textbf{Value (run-1, representative)} \\
\midrule
Lifecycle errors & 0 / 1{,}825 (0\%) \\
\texttt{lifecycle\_op\_ms} median & 1.02\,s \\
\texttt{lifecycle\_op\_ms} p95    & 9.85\,s \\
\bottomrule
\end{tabular}

\medskip\noindent No errors, but Advanced crosses its SLO on the loaded tail.

% ---- 08 ----
\subsection{08 — Multi-User Collaboration}

One host posts; multiple guests poll. Per-app
\texttt{multi\_propagation\_all\_ms}, mean of 5 runs.

\begin{tabular}{@{}lrrrc@{}}
\toprule
\textbf{App} & \textbf{avg} & \textbf{median} & \textbf{p95} & \textbf{threshold} \\
\midrule
Calculator      & 45\,ms & 38\,ms & 38\,ms & \pass \\
Geospatial      & 28\,ms & 22\,ms & 22\,ms & \pass \\
DataExchange    & 40\,ms & 33\,ms & 33\,ms & \pass \\
MonteCarlo      & 37\,ms & 31\,ms & 31\,ms & \pass \\
Advanced        & 57\,ms & 57\,ms & 57\,ms & \pass \\
Analytics       & 38\,ms & 20\,ms & 20\,ms & \pass \\
ClimateAnomaly  & 29\,ms & 26\,ms & 26\,ms & \pass \\
MLTrainer       & 23\,ms & 20\,ms & 20\,ms & \pass \\
\bottomrule
\end{tabular}

\medskip
\noindent All apps pass comfortably. (Note: this test uses a small fixed VU
count per group, so it does not stress the relay the way the ramped tests do —
which is why the heavy apps look fast here.)

% ---- 09 ----
\subsection{09 — WebSocket Presence}

\begin{tabular}{@{}lrrr@{}}
\toprule
\textbf{Metric} & \textbf{median} & \textbf{p95} & \textbf{vs.\ VU load} \\
\midrule
\texttt{ws\_connect\_ms}        & 4\,ms & 6\,ms & flat (no degradation) \\
\texttt{ws\_join\_roundtrip\_ms} & 4\,ms & 6\,ms & flat \\
\bottomrule
\end{tabular}

\medskip\noindent Connect stays $\approx$\,4\,ms from 1 to 50 VUs.

% ------------------------------------------------------------------
\section{Analysis}
% ------------------------------------------------------------------

REST is correct on every reliability dimension: zero data loss, zero
cross-contamination, zero lifecycle errors across all five runs.

On latency, REST is fast at the median for every app (tens of ms on relay and
propagation). The characteristic weakness is the \textbf{loaded tail}: p95
climbs to several seconds because the synchronous Spring Boot relay serialises
requests under concurrency. The POST/POLL split (\S\ref{sec:postpoll}) localises
this to the POST leg — the relay itself — while the Redis read stays cheap.
This is the expected trade-off of a single central relay point.

MLTrainer is a distinct case: its large median (31\,s on relay) is intrinsic R
computation, the same application property seen under Kafka, not a REST relay
cost.

\paragraph{Comparison hook (REST vs.\ Kafka).} Both architectures are equally
correct. The latency \emph{shape} differs: under Kafka the heavy apps pay in
the POLL (waiting for their own compute) while light apps are flat at
$\approx$\,265\,ms; under REST the cost appears in the POST leg as a loaded
tail that affects all apps as concurrency rises. A like-for-like latency-vs-VU
comparison requires the VU tagging fix on tests 01/02/03/08.

\paragraph{Caveats.} (1) Throughput excluded as invalid. (2) Latency tests
aggregate across VU levels (no \texttt{vu\_bucket} tag); medians and p95 should
be read as typical-case and loaded-tail respectively, not as one operating
point.

\end{document}
